\section{Questão 04: Destilação de Conhecimento}
\label{sec:q4}

\subsection{Objetivo e configuração}

A Questão 04 investiga a destilação de conhecimento entre dois modelos da família
Qwen2.5: um professor \texttt{Qwen2.5-3B-Instruct} e um aluno
\texttt{Qwen2.5-1.5B-Instruct}. A escolha de modelos da mesma família é deliberada,
pois compartilham vocabulário e tokenizador, o que reduz incompatibilidades de
distribuição. A técnica empregada é a destilação por \textit{Chain-of-Thought}
(Seção~\ref{sec:fund-destilacao}): o professor produz, para cada trecho, um
raciocínio passo a passo e uma resposta final, e o aluno é ajustado sobre esses
exemplos por meio de LoRA. A Tabela~\ref{tab:q4-config} resume a configuração.

\begin{table}[ht]
  \centering
  \caption{Configuração da destilação da Questão 04.}
  \label{tab:q4-config}
  \footnotesize
  \begin{tabularx}{0.92\textwidth}{@{}l X@{}}
    \toprule
    \textbf{Item} & \textbf{Valor} \\
    \midrule
    Professor / Aluno & Qwen2.5-3B-Instruct ($\sim$6~GB) / Qwen2.5-1.5B-Instruct ($\sim$3~GB) \\
    Técnica & Destilação por \textit{Chain-of-Thought} + LoRA (\texttt{SFTTrainer}) \\
    Dataset & DOMPI-2025, 100 trechos representativos (geração sintética) \\
    Posto $r$ / fator $\alpha$ & 16 / 32, \textit{dropout} 0,05 \\
    Módulos adaptados & \texttt{q\_proj}, \texttt{v\_proj} \\
    Épocas / lote efetivo & 3 / 16 (2 $\times$ 8 de acumulação de gradiente) \\
    Taxa de aprendizado & $2 \times 10^{-4}$, escalonador linear com \textit{warmup} \\
    Comprimento máximo / precisão & 512 tokens / \texttt{fp16} \\
    Ambiente & GPU NVIDIA RTX 4050 (6~GB) \\
    \bottomrule
  \end{tabularx}
\end{table}

\subsection{Dataset sintético}

O dataset de destilação foi gerado a partir de 100 trechos do corpus DOMPI-2025.
Para cada trecho, o professor recebeu um \textit{prompt} solicitando raciocínio
explícito e resposta final em formato JSON, com os campos \texttt{instruction},
\texttt{input}, \texttt{output}, \texttt{reasoning} e \texttt{answer}. As instruções
cobrem extração de informações típicas de diários oficiais, como CNPJ, CEP,
endereços, valores de contrato, números de decreto e modalidades de licitação. O
aluno foi então ajustado sobre esses pares, aprendendo tanto o formato das respostas
quanto o padrão de raciocínio do professor.

\subsection{Protocolo de avaliação}

A qualidade foi medida por dois métodos complementares sobre um \textit{benchmark} de
100 perguntas com respostas de referência. O primeiro é um critério automático de
correspondência textual normalizada (\textit{substring} e sobreposição de palavras
$\geq 70\%$). O segundo é uma avaliação semântica por \textit{LLM-as-Judge}, seguindo
o protocolo de~\cite{zheng2023judging}, em que um LLM juiz classifica cada resposta
como correta ou incorreta, tolerando diferenças de formatação, mas penalizando
contradições e dados inventados. Além disso, as três métricas quantitativas usuais
(entropia cruzada, perplexidade e acurácia de tokens) foram calculadas para o
professor e para o aluno, antes e depois da destilação.

\subsection{Resultados quantitativos}

A Tabela~\ref{tab:q4-quant} mostra que a destilação reduziu a perplexidade do aluno
de 5,49 para 4,23 (queda de 22,9\%) e a entropia cruzada em 15,2\%, elevando a
acurácia de tokens em 2,98 pontos percentuais. O aluno destilado chegou a
\emph{superar o professor} nas três métricas, resultado explicado pelo viés de
ajuste: as métricas são calculadas sobre a mesma distribuição usada no treino do
aluno, que se especializou no domínio, enquanto o professor permanece um modelo
genérico.

\begin{table}[ht]
  \centering
  \caption{Métricas quantitativas antes e depois da destilação (Questão 04).}
  \label{tab:q4-quant}
  \footnotesize
  \begin{tabular}{@{}l c c c@{}}
    \toprule
    \textbf{Modelo} & \textbf{Perplexidade} & \textbf{Entropia cruzada} & \textbf{Acurácia de tokens} \\
    \midrule
    Professor (3B)        & 4,50 & 1,5050 & 68,08\% \\
    Aluno antes (1.5B)    & 5,49 & 1,7022 & 65,65\% \\
    Aluno depois (1.5B)   & \textbf{4,23} & \textbf{1,4434} & \textbf{68,65\%} \\
    \midrule
    Variação do aluno     & $-22{,}9\%$ & $-15{,}2\%$ & $+2{,}98$~pp \\
    \bottomrule
  \end{tabular}
\end{table}

\subsection{Resultados qualitativos: dois métodos divergentes}

O achado mais relevante da questão é a divergência entre os dois métodos de avaliação
qualitativa (Tabela~\ref{tab:q4-qual}). Pela métrica automática de \textit{substring},
o aluno melhorou de 36\% para 42\% de acertos (+6~pp), com 16 questões melhorando
contra 10 regressões. Já pela avaliação semântica por \textit{LLM-as-Judge}, o aluno
\emph{piorou}, de 69\% para 62\% (-7~pp), com 17 regressões contra 10 melhoras; em 52
das 100 perguntas o julgamento não mudou (ambas corretas) e em 21 ambas seguiram
incorretas.

\begin{table}[ht]
  \centering
  \caption{Acertos no \textit{benchmark} de 100 perguntas por método de avaliação.}
  \label{tab:q4-qual}
  \footnotesize
  \begin{tabular}{@{}l c c c@{}}
    \toprule
    \textbf{Método} & \textbf{Aluno antes} & \textbf{Aluno depois} & \textbf{Variação} \\
    \midrule
    \textit{Substring} / sobreposição & 36\% & \textbf{42\%} & $+6$~pp \\
    \textit{LLM-as-Judge} (semântico)  & \textbf{69\%} & 62\% & $-7$~pp \\
    \bottomrule
  \end{tabular}
\end{table}

A explicação é que o ajuste por LoRA aproximou o \emph{formato} das respostas do
aluno ao do gabarito (favorecendo a correspondência textual), mas em vários casos
introduziu alucinações, endereços e números de CNPJ/CEP inexistentes no documento, ou
fez o modelo negar informações que antes acertava. A métrica de \textit{substring}
não penaliza esses erros quando o trecho correto também aparece; o juiz semântico,
sim.

\subsection{Houve transferência de conhecimento?}

Sim, houve transferência, mas com evidências mistas quanto à sua qualidade. As três
métricas quantitativas confirmam a transferência de forma inequívoca (perplexidade e
entropia cruzada menores, acurácia maior), e o aluno chegou a superar o professor no
domínio avaliado. Contudo, ao nível da corretude semântica das respostas, o aluno
piorou. A lição central é que menor perplexidade não implica respostas mais confiáveis
do ponto de vista factual: ganhos de formato e concisão podem mascarar perda de
confiabilidade que só a avaliação semântica, ou humana, é capaz de capturar.
